\documentclass[conference]{IEEEtran}
\IEEEoverridecommandlockouts

\usepackage{cite}
\usepackage{amsmath,amssymb,amsfonts}
\usepackage{algorithmic}
\usepackage{graphicx}
\usepackage{textcomp}
\usepackage{xcolor}
\usepackage{float}
\setlength{\textfloatsep}{7pt plus 1pt minus 2pt}
\setlength{\floatsep}{7pt plus 1pt minus 2pt}
\setlength{\intextsep}{7pt plus 1pt minus 2pt}
\setlength{\abovecaptionskip}{3pt plus 1pt minus 1pt}
\setlength{\belowcaptionskip}{1pt plus 1pt minus 1pt}
\setlength{\tabcolsep}{4pt}
\newcommand{\placeholderfigure}[1]{%
\IfFileExists{#1}{%
\includegraphics[width=\linewidth]{#1}%
}{%
\fbox{\parbox[c][1.2in][c]{0.92\linewidth}{\centering\scriptsize Placeholder figure:\par\texttt{\detokenize{#1}}}}%
}%
}
\def\BibTeX{{\rm B\kern-.05em{\sc i\kern-.025em b}\kern-.08em
    T\kern-.1667em\lower.7ex\hbox{E}\kern-.125emX}}
\begin{document}
\raggedbottom

\title{Experimental Characterization of the Output Stage of a Functional Electrical Stimulator Based on a Howland Current Source}

\author{
\IEEEauthorblockN{T. P. Silva}
\IEEEauthorblockA{
\textit{Laboratório de Engenharia Biomédica} \\
\textit{Centro Universitário FEI} \\
São Bernardo do Campo, Brasil \\
tpsilva@fei.edu.br}
}
\maketitle

\begin{abstract}
Functional Electrical Stimulation (FES) systems require predictable current control to ensure repeatable clinical outcomes and stable embedded control. This paper characterizes the relationship between digital-to-analog converter (DAC) control voltage and output current in a Howland current source output stage based on the STIMGRASP FES architecture. As this architecture lacks real-time current feedback, the output stage operates in an open-loop configuration. Therefore, for experimental characterization only, stimulation current was measured externally using a shunt resistor. A dedicated firmware was developed to generate DAC ramps under multiple resistive load conditions. After excluding compliance-limited samples, linearity and regression analyses were performed. Pearson and Spearman correlation coefficients were $<-0.999$ for all loads, confirming highly linear and monotonic behavior in the selected operating region. A single global linear model was identified for use as an embedded DAC estimation model. Finally, this work quantifies how load-dependent compliance limits may compromise predictability in the absence of feedback, providing a characterization that complements the original STIMGRASP investigation.
\end{abstract}


\begin{IEEEkeywords}
Functional Electrical Stimulation, Howland current source, linearity analysis, open-loop control.
\end{IEEEkeywords}

\section{Introduction}

Functional Electrical Stimulation (FES) systems electrically activate peripheral nerves or muscles in order to assist or restore motor functions in rehabilitation and assistive applications \cite{FES_REVIEW_1}, \cite{FES_REVIEW_2}. In current-controlled stimulation, the delivered current amplitude is a central parameter because it is directly associated with the recruited neuromuscular response, user comfort, and repeatability of the stimulation protocol \cite{CURRENT_CONTROLLED_STIMULATION}. For this reason, the relationship between the embedded command variable and the actual output current must be experimentally characterized.

This work is related to the STIMGRASP FES master's thesis developed at Centro Universitário FEI by Barelli, which proposed a portable neuromuscular electrical stimulator for grasp restoration \cite{STIMGRASP_THESIS}. The original work describes a digitally controlled stimulation architecture based on a modified Howland current source and discusses linear stimulation ramps and current-controlled stimulation behavior. In that investigation, predictable current control behavior is assumed as part of the stimulation architecture. However, detailed DAC-versus-current characterization curves under different load conditions are not presented, and the behavior implied by the architecture is not quantitatively characterized through correlation, regression, compliance, or load-consistency analyses.

The present work extends the original investigation by experimentally evaluating the output-stage transfer characteristic of the FES circuit. The purpose is not to redesign the STIMGRASP hardware or to provide a clinical validation of the stimulator. Instead, this study focuses on output-stage characterization: the relationship between DAC control voltage and delivered current, the effect of different resistive loads, the identification of a valid compliance region, and the derivation of an embedded DAC estimation model for open-loop operation.

The STIMGRASP circuit does not include current feedback to the microcontroller unit (MCU). Consequently, the output stage operates as an open-loop stimulation system, where the delivered current is not directly monitored during operation. Under this condition, predictable stimulation behavior depends on a sufficiently linear and repeatable relationship between DAC control voltage and output current within the valid operating region. In this work, the experimentally identified DAC-to-current transfer model is used for embedded DAC estimation. The user or application configures a desired stimulation current, and the firmware estimates the DAC voltage required to produce it. If the skin-electrode impedance increases, electrode contact becomes poor, or the load changes substantially, the output stage may reach voltage compliance and the delivered current may be lower than the modeled current. Without current feedback, this condition cannot be detected by the MCU in real time.

The main contributions of this paper are: (i) quantitative characterization of the DAC-voltage-to-output-current relationship; (ii) experimental compliance-region analysis based on load voltage; (iii) statistical validation of linearity using Pearson and Spearman correlations and regression metrics; (iv) regression analysis by load and a single global model; and (v) discussion of the practical limitations of open-loop current estimation in the absence of closed-loop current monitoring.

\section{Background}

\subsection{Functional Electrical Stimulation}

FES uses electrical pulses to evoke muscle contractions for functional assistance or rehabilitation. Stimulation systems typically allow the configuration of pulse amplitude, width, frequency, and waveform shape. Among these variables, amplitude control is particularly important because it modulates the intensity of nerve and muscle recruitment \cite{FES_REVIEW_1}, \cite{CURRENT_CONTROLLED_STIMULATION}.

The present study is restricted to output-stage characterization under controlled laboratory conditions. Resistive loads were used to represent controlled variations in equivalent human impedance and electrode contact quality while preserving repeatability. Such loads are useful for instrumentation analysis, but they do not fully reproduce the nonlinear, time-varying, and electrode-dependent behavior of the skin-electrode interface \cite{ELECTRODE_SKIN_IMPEDANCE}.

\subsection{Howland Current Source and Compliance Voltage}

The Howland current source is commonly used when a voltage command must be converted into a load current \cite{HOWLAND_CURRENT_SOURCE}. Ideally, the output current is determined by the input control voltage and resistor network, and it remains independent of the load impedance. In practice, load-independent behavior is maintained only while the amplifier output can provide the voltage required by the load. This constraint is described by the compliance voltage \cite{COMPLIANCE_VOLTAGE}.

For a resistive load, the load voltage can be reconstructed from the measured current as
\begin{equation}
V_{\mathrm{load}} = I_{\mathrm{out}} R_{\mathrm{load}},
\label{eq:vload}
\end{equation}
where $I_{\mathrm{out}}$ is expressed in amperes and $R_{\mathrm{load}}$ in ohms. Current-source operation is expected to remain valid while
\begin{equation}
|V_{\mathrm{load}}| < V_{\mathrm{compliance}}.
\label{eq:compliance_condition}
\end{equation}
Samples near or beyond the available output-voltage capability may exhibit compliance-limited behavior, particularly for higher load resistances.

\subsection{Open-Loop DAC Estimation}

Because the evaluated FES circuit does not provide real-time current feedback to the MCU, the embedded controller cannot directly regulate the delivered current during stimulation. Instead, current configuration must rely on an experimentally identified DAC-to-current relationship to estimate the DAC voltage required for a desired stimulation current. The experimentally identified direct model is
\begin{equation}
I_{\mathrm{out}} = a V_{\mathrm{DAC}} + b,
\label{eq:direct_model}
\end{equation}
where $V_{\mathrm{DAC}}$ is the DAC control voltage and $a$ and $b$ are regression coefficients. In the evaluated circuit, $a$ is negative because increasing the DAC voltage decreases the output current.

The embedded DAC estimation equation is obtained by rearranging \eqref{eq:direct_model}:
\begin{equation}
V_{\mathrm{DAC}} = \frac{I_{\mathrm{target}} - b}{a}.
\label{eq:dac_estimation_model}
\end{equation}
Equation \eqref{eq:dac_estimation_model} is valid only within the experimentally selected compliance region. Outside this region, the output current is constrained by available output voltage and load resistance, and the DAC-to-current model may predict a current that cannot be physically delivered once the output stage reaches voltage compliance.

\section{Materials and Methods}

\subsection{Experimental Setup}

The evaluated output stage was an FES circuit based on a Howland current source. The MCU controlled the stimulation command through a DAC voltage input. A dedicated experimental firmware was developed for this characterization. Its only function was to generate a controlled DAC ramp with an initial and final resting interval; it was not the conventional stimulation firmware and did not generate a complete FES pulse train with therapeutic stimulation parameters. The nominal resting or reference point of the DAC command was approximately 1.65 V. This value is not arbitrary: in the STIMGRASP output-stage circuit, it corresponds to the zero-current operating point, for which the output current is approximately 0 mA and the load voltage is consequently approximately 0 V. In this work, the intervals around this 1.65 V reference are referred to as the zero-current resting regime. The DAC voltage was swept from approximately 0 V to 3.3 V.

Figure~\ref{fig:output_stage_schematic} shows the evaluated DAC and Howland-current-source output stage as extracted from the STIMGRASP documentation. The schematic is used only to identify the characterized circuit block and does not imply a hardware redesign.

\begin{figure}[htbp]
\centerline{\placeholderfigure{images/fig_output_stage_schematic.png}}
\caption{Evaluated DAC and Howland-current-source output stage.}
\label{fig:output_stage_schematic}
\end{figure}

Output current was measured experimentally using an external measurement setup with a 10 $\Omega$ shunt resistor. Three resistive loads were evaluated: 1 k$\Omega$, 2 k$\Omega$, and 4.7 k$\Omega$. These loads were selected to emulate controlled variations in equivalent body impedance and electrode contact condition, from lower impedance or better contact to higher impedance or poorer contact, while maintaining repeatable bench-test conditions.

\subsection{Data Acquisition}

Raw data were acquired as CSV files exported by a Tektronix MSO46B oscilloscope using two analog channels. One channel measured the DAC voltage and the other measured the shunt voltage. The oscilloscope configuration is summarized in Table~\ref{tab:oscilloscope}. 

\begin{table}[H]
\caption{Oscilloscope Acquisition Configuration}
\begin{center}
\begin{tabular}{|c|c|}
\hline
\textbf{Parameter} & \textbf{Value} \\
\hline
Oscilloscope model & MSO46B \\
\hline
Channels & CH1 and CH2, analog \\
\hline
Sample interval & 80 $\mu$s \\
\hline
Sample frequency & 12.5 kHz \\
\hline
Vertical unit & V \\
\hline
Horizontal unit & s \\
\hline
\end{tabular}
\label{tab:oscilloscope}
\end{center}
\end{table}

Figure~\ref{fig:raw_dac_shunt} shows a representative raw acquisition for the 1 k$\Omega$ load before any processing. This example is included only to illustrate the format of the acquired oscilloscope data and the preliminary relationship between DAC voltage and shunt voltage. The initial and final resting intervals are preserved in the raw CSV and are visible around the nominal 1.65 V DAC reference level, corresponding to the zero-current operating point of the evaluated output stage. The useful experimental interval corresponds to the DAC ramp.

\begin{figure}[htbp]
\centerline{\placeholderfigure{images/fig_raw_dac_shunt_rest.png}}
\caption{Raw DAC and shunt voltages for the 1 k$\Omega$ load.}
\label{fig:raw_dac_shunt}
\end{figure}

The useful ramp region was extracted from each acquisition before model identification. In this context, ramp extraction means removing the initial and final zero-current resting-regime samples around the 1.65 V DAC reference and preserving only the swept DAC interval. The measured shunt voltage was converted into output current using the 10 $\Omega$ shunt value. The output voltage was not directly measured; instead, it was inferred from \eqref{eq:vload} using the calculated current and nominal load resistance.

\subsection{Data Processing}

The analysis pipeline first extracted the useful DAC-ramp interval corresponding to the experimental sweep from approximately 0 V to 3.3 V, removing the initial and final resting-regime samples around the nominal 1.65 V reference level. The shunt-voltage measurement was then converted into output current in milliamperes. The current signals were grouped by load resistance, as illustrated in Fig.~\ref{fig:current_per_load_time}. This figure shows current evolution over acquisition samples after ramp extraction; it is not a direct current-versus-DAC-voltage representation. Converting the shunt-voltage measurements into output current made the compliance-limited behavior visible in the processed data. In this region, the load current no longer increases proportionally with the DAC-voltage command, because the output stage is constrained by the voltage required across the load.

\begin{figure}[htbp]
\centerline{\placeholderfigure{images/fig_current_per_load_time.png}}
\caption{Calculated output current over acquisition samples.}
\label{fig:current_per_load_time}
\end{figure}

In the next step, the DAC voltage axis was binned into 10 mV intervals to reduce granularity and measurement noise before compliance-region selection and model identification. For each load resistance and each DAC-voltage bin, the representative current value was obtained as the arithmetic mean of all current samples within that interval. The complete binned representation, shown in Fig.~\ref{fig:binned_with_rest}, includes both the linear current-source region and the compliance-limited region. It is therefore appropriate for characterizing full-range behavior and operating limits, but it is not used directly for deriving the linear DAC estimation model.


\begin{figure}[H]
\centerline{\placeholderfigure{images/fig_binned_current_with_rest.png}}
\caption{Binned output current versus DAC voltage.}
\label{fig:binned_with_rest}
\end{figure}

\subsection{Compliance-Limited Region Selection}

The compliance-limited region was identified experimentally from the reconstructed load voltage. For each load, the minimum and maximum measured current values were used to estimate the corresponding minimum and maximum load voltages. The common valid experimental voltage interval was then selected as
\begin{equation}
V_{\min,\mathrm{limit}} = \max(V_{\min,1k}, V_{\min,2k}, V_{\min,4.7k}),
\label{eq:vmin_limit}
\end{equation}
and
\begin{equation}
V_{\max,\mathrm{limit}} = \min(V_{\max,1k}, V_{\max,2k}, V_{\max,4.7k}).
\label{eq:vmax_limit}
\end{equation}
This method avoids hardcoded voltage thresholds and selects a common interval in which all three loads remain experimentally inside the valid region.

The inferred voltage limits are summarized in Table~\ref{tab:voltage_limits}. The resulting common selected voltage interval was approximately -43.803 V to 46.330 V. Samples whose inferred $V_{\mathrm{load}}$ remained inside this interval formed the valid linear operating dataset. This distinction is essential: the complete binned data describe the full behavior of the circuit, whereas only the selected valid operating data are used for linearity, regression, residual, load-consistency, and embedded DAC estimation analyses.

\begin{table}[H]
\caption{Experimental Load-Voltage Limits Inferred From Measured Current}
\begin{center}
\begin{tabular}{|c|c|c|}
\hline
\textbf{Load} & \textbf{$V_{\min}$ (V)} & \textbf{$V_{\max}$ (V)} \\
\hline
1 k$\Omega$ & -43.803 & 46.330 \\
\hline
2 k$\Omega$ & -46.225 & 49.560 \\
\hline
4.7 k$\Omega$ & -47.619 & 47.776 \\
\hline
Common interval & -43.803 & 46.330 \\
\hline
\end{tabular}
\label{tab:voltage_limits}
\end{center}
\end{table}

After removing samples outside this common compliance interval, Fig.~\ref{fig:full_current_vs_dac} shows the output current versus DAC voltage for the three evaluated loads. The remaining curves correspond to the selected linear operating region used for model identification. Their close agreement indicates that, once compliance-limited samples are excluded, the output current is mainly determined by the DAC command rather than by the load resistance.

\begin{figure}[H]
\centerline{\placeholderfigure{images/fig_full_current_vs_dac.png}}
\caption{Output current versus DAC voltage after compliance-limited sample removal.}
\label{fig:full_current_vs_dac}
\end{figure} 


\subsection{Statistical Analysis and Model Identification}

Pearson and Spearman correlations were computed for the selected linear operating region. Pearson correlation was used to quantify linear association, while Spearman correlation was used to evaluate monotonic association \cite{PEARSON_SPEARMAN_COMPARISON}, \cite{CORRELATION_COEFFICIENTS_INTERPRETATION}. Compliance-limited samples were excluded before statistical analysis, so the reported linearity metrics refer only to the hardware-supported operating region used for model identification.

After correlation analysis, the DAC-to-current relationship was modeled as a linear regression problem in the form of \eqref{eq:direct_model}. Linear regression was first performed independently for each load to quantify load-specific slope, intercept, coefficient of determination, and estimation error in the valid region.

A second regression was then performed globally using all selected samples from the 1 k$\Omega$, 2 k$\Omega$, and 4.7 k$\Omega$ loads. For this global fit, each point was represented by a pair $(V_{\mathrm{DAC}}, I_{\mathrm{out}})$, with the load label retained only as auxiliary information. All points from the three loads were combined into a single dataset and used simultaneously in one linear regression. The objective was to identify a single embedded DAC estimation model valid across the evaluated loads inside the hardware-supported compliance interval.

The global regression was used as the basis for the open-loop DAC estimation model. Residuals were analyzed afterward to assess systematic deviations from this global model.

\newpage
\section{Results}

\subsection{Linear Operating Region}

After removing samples outside the common valid voltage interval, the curves for 1 k$\Omega$, 2 k$\Omega$, and 4.7 k$\Omega$ overlap closely in the selected region. This indicates that the output current is primarily determined by the DAC voltage and not by the load resistance. The correlation results for this region are summarized in Table~\ref{tab:correlation}. Pearson and Spearman correlation coefficients were very close to -1 for all loads, confirming highly linear and monotonic behavior after compliance-limited samples were removed. The negative sign is expected because, in the evaluated circuit, increasing DAC voltage decreases the output current.

\begin{table}[!htbp]
\caption{Pearson and Spearman Correlation Results in the Linear Operating Region}
\begin{center}
\scriptsize
\renewcommand{\arraystretch}{0.92}
\begin{tabular}{|c|c|c|c|c|c|}
\hline
\textbf{Load} & \textbf{$r$} & \textbf{$p_r$} & \textbf{$\rho$} & \textbf{$p_\rho$} & \textbf{$n$} \\
\hline
1 k$\Omega$ & -0.999909 & 0.000000 & -1.000000 & 0.000000 & 323 \\
\hline
2 k$\Omega$ & -0.999899 & 0.000000 & -1.000000 & 0.000000 & 164 \\
\hline
4.7 k$\Omega$ & -0.999824 & 0.000000 & -1.000000 & 0.000000 & 74 \\
\hline
\end{tabular}
\renewcommand{\arraystretch}{1}
\vspace{1mm}

\parbox{\columnwidth}{\scriptsize $r$: Pearson correlation coefficient; $p_r$: Pearson p-value; $\rho$: Spearman correlation coefficient; $p_\rho$: Spearman p-value; $n$: number of samples.}
\label{tab:correlation}
\end{center}
\end{table}

\subsection{Linear Regression and Global Model}

Linear regression in the selected operating region showed high linearity for all loads. Table~\ref{tab:load_regression} summarizes the per-load regression coefficients, while Table~\ref{tab:load_error_metrics} reports the corresponding estimation errors.

\begin{table}[!htbp]
\caption{Linear Regression Coefficients by Load in the Valid Region}
\begin{center}
\scriptsize
\renewcommand{\arraystretch}{0.92}
\begin{tabular}{|c|c|c|c|c|}
\hline
\textbf{Load} & \textbf{$a$ (mA/V)} & \textbf{$b$ (mA)} & \textbf{$R^2$} & \textbf{SE (mA)} \\
\hline
1 k$\Omega$ & -28.1153 & 46.6128 & 0.999819 & 0.02114 \\
\hline
2 k$\Omega$ & -27.6292 & 45.5749 & 0.999799 & 0.03079 \\
\hline
4.7 k$\Omega$ & -26.2052 & 42.9934 & 0.999648 & 0.05797 \\
\hline
\end{tabular}
\renewcommand{\arraystretch}{1}
\vspace{1mm}

\parbox{\columnwidth}{\scriptsize $a$: fitted slope; $b$: fitted intercept; $R^2$: coefficient of determination; SE: standard error of the regression.}
\label{tab:load_regression}
\end{center}
\end{table}

\begin{table}[!htbp]
\caption{Per-Load Error Metrics in the Valid Region}
\begin{center}
\scriptsize
\renewcommand{\arraystretch}{0.92}
\begin{tabular}{|c|c|c|c|}
\hline
\textbf{Load} & \textbf{MAE (mA)} & \textbf{RMSE (mA)} & \textbf{Max. (mA)} \\
\hline
1 k$\Omega$ & 0.2878 & 0.3532 & 0.7411 \\
\hline
2 k$\Omega$ & 0.1535 & 0.1855 & 0.4206 \\
\hline
4.7 k$\Omega$ & 0.0918 & 0.1051 & 0.2611 \\
\hline
\end{tabular}
\renewcommand{\arraystretch}{1}
\vspace{1mm}

\parbox{\columnwidth}{\scriptsize MAE: mean absolute error; RMSE: root mean square error; Max.: maximum absolute error.}
\label{tab:load_error_metrics}
\end{center}
\end{table}

The preferred embedded model is a single global linear model identified using all three loads within the valid linear region. Figure~\ref{fig:global_model} shows the experimental curves and the resulting global linear model. The close overlap indicates that the open-loop DAC estimation model can represent the evaluated output stage across the tested loads, provided that the load voltage remains inside the compliance interval.

\begin{figure}[H]
\centerline{\placeholderfigure{images/fig_global_linear_model.png}}
\caption{Experimental current curves and global linear model.}
\label{fig:global_model}
\end{figure}

\begin{table}[!htbp]
\caption{Global Linear Model Performance Metrics}
\begin{center}
\scriptsize
\renewcommand{\arraystretch}{0.92}
\begin{tabular}{|l|c|}
\hline
\textbf{Metric} & \textbf{Value} \\
\hline
Slope $a$ (mA/V) & -28.040197 \\
\hline
Intercept $b$ (mA) & 46.353724 \\
\hline
$R^2$ & 0.999695 \\
\hline
MAE (mA) & 0.328845 \\
\hline
RMSE (mA) & 0.385635 \\
\hline
Max. error (mA) & 0.930634 \\
\hline
Samples & 561 \\
\hline
\end{tabular}
\renewcommand{\arraystretch}{1}
\vspace{1mm}

\parbox{\columnwidth}{\scriptsize $a$: fitted slope; $b$: fitted intercept; $R^2$: coefficient of determination; MAE: mean absolute error; RMSE: root mean square error.}
\label{tab:global_model}
\end{center}
\end{table}

\subsection{Residuals and Load Consistency}

Figure~\ref{fig:residuals} shows the residuals of the global linear model. The residual distribution was used to identify possible systematic deviations as a function of DAC voltage. Such deviations are important for embedded implementation because they define the expected open-loop estimation error in the absence of real-time current feedback.

\begin{figure}[!htbp]
\centerline{\placeholderfigure{images/fig_residuals.png}}
\caption{Residuals of the global linear model.}
\label{fig:residuals}
\end{figure}

\section{Discussion}

The results indicate that the evaluated Howland-based FES output stage exhibits strong linearity between DAC control voltage and output current within the experimentally selected valid region. Pearson correlation coefficients close to -1, Spearman correlations equal or close to -1, and $R^2$ values above 0.9996 demonstrate that the relationship is highly linear and monotonic after compliance-limited samples are removed. The negative sign of the correlations and fitted slopes is consistent with the polarity of the evaluated circuit, in which increasing DAC voltage decreases output current.

The nonlinear appearance of the complete DAC sweep is explained by compliance-limited behavior, not by a general loss of current-source linearity. The complete dataset is useful because it characterizes the experimental operating limits of the output stage. However, it mixes the linear current-source regime with a region where output current is constrained by the available voltage across the load. For this reason, compliance-limited samples were excluded from the regression analysis and from the derivation of the embedded DAC estimation model.

The global linear model is practically useful for open-loop firmware current configuration. With coefficients $a$ and $b$ identified from the selected valid operating data, the MCU can estimate the required DAC voltage from the desired stimulation current using \eqref{eq:dac_estimation_model}. This supports the open-loop control strategy implied by the original STIMGRASP architecture while adding quantitative calibration and validation to the DAC-to-current transfer behavior described conceptually in the dissertation \cite{STIMGRASP_THESIS}.

At the same time, the open-loop DAC estimation model has an important limitation: it is valid only inside the compliance region. In real use, increased skin-electrode impedance, poor electrode contact, or significant load changes may require an output voltage beyond the available compliance. In that case, the real delivered current may be lower than the modeled current. Because the evaluated system does not include current feedback to the MCU, the firmware cannot detect this condition in real time. This limitation is directly related to the distinction between open-loop and closed-loop stimulation control discussed in the broader context of current-controlled FES systems.

Another limitation is the use of resistive loads. Although resistors provide repeatable and well-defined test conditions for output-stage characterization, they do not reproduce the nonlinear and time-varying impedance of the skin-electrode interface. Future work should therefore include tests with more realistic loads or electrode models. Hardware or firmware extensions for current feedback, compliance monitoring, or closed-loop current regulation should also be considered to improve robustness under changing load conditions.

\section{Conclusion}

This paper presented an experimental characterization of the output-stage linearity of an FES circuit based on a Howland current source and related to the STIMGRASP stimulation architecture. A DAC ramp was applied over approximately 0 V to 3.3 V, and output current was measured under 1 k$\Omega$, 2 k$\Omega$, and 4.7 k$\Omega$ resistive loads. The complete dataset was used to characterize compliance-limited behavior, while a common valid load-voltage interval of approximately -43.80 V to 46.33 V was selected for linearity and regression analyses.

Within this valid region, the DAC-to-current relationship was highly linear and monotonic for all tested loads, with Pearson and Spearman correlation coefficients close to -1. A single global linear model using all three loads is therefore appropriate for embedded open-loop DAC estimation, provided that operation remains within the compliance region. The main practical limitation is that, without real-time current feedback, the MCU cannot detect when load changes or skin-electrode impedance cause compliance-limited operation. The study complements the original STIMGRASP investigation by adding quantitative transfer characterization, statistical validation, compliance-region analysis, and open-loop DAC estimation discussion.

\bibliographystyle{IEEEtran}
\bibliography{references}

\end{document}
